%; whizzy paragraph -pdf xpdf -latex ./whizzypdfptex.sh
%; whizzy-paragraph "^\\\\begin{frame}"
% latex beamer presentation.
% platex, latex-beamer でコンパイルすることを想定。 

%     Tokyo Debian Meeting resources
%     Copyright (C) 2011 Junichi Uekawa
%     Copyright (C) 2011 Nobuhiro Iwamatsu

%     This program is free software; you can redistribute it and/or modify
%     it under the terms of the GNU General Public License as published by
%     the Free Software Foundation; either version 2 of the License, or
%     (at your option) any later version.

%     This program is distributed in the hope that it will be useful,
%     but WITHOUT ANY WARRANTY; without even the implied warreanty of
%     MERCHANTABILITY or FITNESS FOR A PARTICULAR PURPOSE.  See the
%     GNU General Public License for more details.

%     You should have received a copy of the GNU General Public License
%     along with this program; if not, write to the Free Software
%     Foundation, Inc., 51 Franklin St, Fifth Floor, Boston, MA  02110-1301 USA

\documentclass[cjk,dvipdfmx,12pt]{beamer}
\usetheme{Tokyo}
\usepackage{monthlypresentation}

%  preview (shell-command (concat "evince " (replace-regexp-in-string "tex$" "pdf"(buffer-file-name)) "&")) 
%  presentation (shell-command (concat "xpdf -fullscreen " (replace-regexp-in-string "tex$" "pdf"(buffer-file-name)) "&"))
%  presentation (shell-command (concat "evince " (replace-regexp-in-string "tex$" "pdf"(buffer-file-name)) "&"))

%http://www.naney.org/diki/dk/hyperref.html
%日本語EUC系環境の時
\AtBeginDvi{\special{pdf:tounicode EUC-UCS2}}
%シフトJIS系環境の時
%\AtBeginDvi{\special{pdf:tounicode 90ms-RKSJ-UCS2}}

\title{Debian とはなにか？}
\subtitle{第81回 2011年10月度 勉強会}
\author{岩松 信洋 iwamatsu@debian.org\\IRC nick: iwamatsu}
\date{2011年10月22日}
\logo{\includegraphics[width=8cm]{image200607/openlogo-light.eps}}

\begin{document}

\frame{\titlepage{}}


\begin{frame}{自己紹介}
  \begin{itemize}
  \item 岩松です。
  \item Debian Developer です。
  \item Bluetooth周り とか opencvとかメンテナンスしています。
  \item Renesas SH4 ポーターです。
  \item 普段はカーネルの開発とか、ユーザランド整備とかしてます。
  \end{itemize}
\end{frame}


\begin{frame}{アジェンダ}
  \begin{itemize}
  \item Debian とは？  / Debian の概要
  \item 特徴 / Debian の特徴
  \item ユーザはどうしたらいいのか？
  \end{itemize}
\end{frame}


\begin{frame}{Debianとは？}

\begin{itemize}
\item Debian Project の略称、またはDebian OSそのものを指す
\item フリーかつオープンなOSを作る完全ボランティアベースのプロジェクト
\item 歴史が長くいい意味で保守的なLinuxディストリビューションの一つ
\item 公式開発者は約1000名。
\item 非公式な開発者やパッケージメンテナ、翻訳者などを入れると5000名以上
\item 日本の開発者は約40名。
\item Debian JP Project
日本でのDebian環境のサポート、開発者育成、ユーザサポートなどを行う団体。
\end{itemize}
\end{frame}

\begin{frame}{特徴}
\begin{itemize}
\item フリーなソフトウェアで構成されている
\item オープンな開発をしている
\item バイナリベースのディストリビューション
\item 豊富なパッケージ数
\item ポリシーに基づいたパッケージ
\item 強力なパッケージングシステム
\item アップデートが容易
\item 豊富なサポートCPUアーキテクチャ
\item Linux以外のカーネルもサポートする
\item 他のOSSプロジェクトと関連が強い
\item 派生しているディストリビューションの多さ
\end{itemize}
\end{frame}


\begin{frame}{フリーなソフトウェアで構成されている}
Debian フリーソフトウェアガイドライン (DFSG) 

\begin{itemize}
\item 何台のマシンにもソフトウェアをインストールしても良い。
\item 何人の人がソフトウェアを同時に使用しても良い。
\item 何個でもソフトウェアのコピーを作ってもいいし、それを誰にあげても構わない。(フリーもしくはオープンな再配布)
\item ソフトウェアの改変に対する制約が無い。(特定の通告を変えない事を除く)
\item そのソフトウェアを配布や売る事に対する制約が無い。
\end{itemize}
\end{frame}

\begin{frame}{オープンな開発をしている}

\begin{itemize}
\item 完全にボランティアベースの団体
\item 特定の企業によって方針を左右されない
\item なにか提案したいときは、選挙で決めることができる。
\item 議論の経過なども全て公開
\item バグのオープンからクローズまでも全て公開
\end{itemize}

\end{frame}


\begin{frame}{バイナリベースのディストリビューション}

\begin{itemize}
\item バイナリベース / Debian, Ubuntu, Fedora...
\begin{itemize}
\item 導入が容易
\item 問題の共有が容易
\item エコ
\end{itemize}
\item ソース ベース Gentoo, etc
\begin{itemize}
\item 自分でコンパイル。環境に合わせた最適化がしやすい。
\item 問題の再現性が低い（俺のところでは再現しませんが、なにか？とか）
\item エコじゃない
\end{itemize}
\end{itemize}
\end{frame}

\begin{frame}{ポリシーに基づいたパッケージ}
\begin{itemize}
\item Debian に要求されるいくつかの必要条件について Debian Policyと
して定められている。
\item ポリシーには、Debian アーカイブの構成と内容、
オペレーティングシステムとしての Debian の設計に
関しても定めらている。
\item それぞれのパッケージがディストリビューションに受け入れら
れるために満たさなければならない技術的な必要条件も含まれる。
\item パッケージはこのポリシーに従って作成されるので、どこかの
パッケージングシステムのようにディストリビューションによって
仕様が異なることがない。
\end{itemize}
\end{frame}


\begin{frame}{豊富なパッケージ数}
\begin{itemize}
\item 約30000パッケージ。
\item 毎月20個ぐらい増えている。
\item Gentoo で約15000パッケージ。
\item Ubuntu はもっと少ない（mainだけだし）。
\end{itemize}
\end{frame}


\begin{frame}{強力なパッケージングシステム}
\begin{itemize}
\item dpkg
\item Debian Policy
\item パッケージ間の依存関係
\end{itemize}

\end{frame}


\begin{frame}{アップデートが容易}
\begin{itemize}

\item APT/aptitude 様々。
\item 現在の使っているバージョンから順にアップデートできる。
\item 再インストール？なんすかそれ？
\end{itemize}

\end{frame}


\begin{frame}{豊富なサポートCPUアーキテクチャ}
\begin{itemize}
\item 現時点で正式サポートCPUアーキテクチャは 11
\item 次期リリースに向けてサポート準備中は 10
\item Gentooはもっと多いけど、動くかわからん....。
\end{itemize}

\begin{minipage}[t]{0.5\hsize}

現在サポートしているアーキテクチャ。
\begin{itemize}
  \item amd64
  \item armel
  \item hurd-i386
  \item i386
  \item ia64
  \item kfreebsd-amd64
  \item kfreebsd-i386
  \item mips
  \item mipsel
  \item powerpc
  \item s390
  \item sparc
\end{itemize}

\end{minipage}
\begin{minipage}[t]{0.5\hsize}

サポート予定のアーキテクチャ
\begin{itemize}
  \item alpha
  \item armhf
  \item avr32
  \item hppa
  \item m68k
  \item powerpcspe
  \item s390x
  \item sh4
  \item sparc64
\end{itemize}

\end{minipage}

\end{frame}



\begin{frame}{Linux以外のカーネルもサポートする}
\begin{itemize}
\item サポートはLinuxカーネルだけではない。
\item FreeBSDのカーネル \\
Debian GNU/kFreeBSD 
\item GNU Hurd \\
Debian GNU/Hurd
\item Minix, NetBSD カーネル 向けも開発中。

\end{itemize}

\end{frame}

\begin{frame}{他のOSSプロジェクトと関連が強い}
\begin{itemize}
\item 開発元をやっている人がけっこう多い。
\item パッチが取り込まれやすい。
\end{itemize}
\end{frame}


\begin{frame}{派生しているディストリビューションの多さ}

\begin{itemize}
\item 派生ディストリビューション数 １位\\
129以上。UbuntuやKnoppix、Vyatta etc....
\item 2番目に多いのはFedoraベースの 63
\item 派生として分散させているだけでなく、派生したディストリビューションの
成果を本家であるDebianに取り組む仕組みもある。
\item 派生が多い理由は、いままで説明した特徴によるところが大きい。
\end{itemize}
\end{frame}

\begin{frame}{まとめ}
\begin{itemize}
\item フリーなソフトウェアで構成されている
\item オープンな開発をしている
\item バイナリベースのディストリビューション
\item 豊富なパッケージ数
\item ポリシーに基づいたパッケージ
\item 強力なパッケージングシステム
\item アップデートが容易
\item 豊富なサポートCPUアーキテクチャ
\item Linux以外のカーネルもサポートする
\item 他のOSSプロジェクトと関連が強い
\item 派生しているディストリビューションの多さ
\end{itemize}
\end{frame}


\begin{frame}{どういう風に使えばいいの？}
\begin{itemize}[<+->]
\item 開発に使いたいなら、Debian。\\
\item デスクトップやノートPCで使いたいなら Debian。
\item サーバで使いたいなら、Debian。\\
\item ようするにDebian。
\item Debianには使う楽しさと、創り上げる楽しさの両方がある。
\end{itemize}


\end{frame}


\begin{frame}{最後に}
\begin{itemize}
\item ぜひDebianを使って、フィードバックをください。
\item そして開発に興味がある人は開発に参加してみてください。手取り足取り教えます。
\item みんなでDebianを良いものにしていきましょう。
\end{itemize}
\end{frame}

\end{document}

;;; Local Variables: ***
;;; outline-regexp: "\\([ 	]*\\\\\\(documentstyle\\|documentclass\\|emtext\\|section\\|begin{frame}\\)\\*?[ 	]*[[{]\\|[]+\\)" ***
;;; End: ***
